\documentclass[twocolumn]{article}
\usepackage[utf8]{inputenc}
\usepackage[T1]{fontenc}
\usepackage[spanish, es-tabla]{babel} % es-tabla cambia "Cuadro" por "Tabla"
\usepackage{amsmath, amssymb}
\usepackage{algorithm}
\usepackage{algpseudocode}
\usepackage{booktabs}
\usepackage{geometry}
\usepackage{hyperref}

\geometry{margin=1.5cm, top=1.8cm, bottom=1.8cm}
\hypersetup{colorlinks=true, urlcolor=blue}

\title{\textbf{Prefix Sum: De $O(n)$ a $O(1)$}\\
\large Precomputación como estrategia de escalabilidad masiva}
\author{Ashly Hernandez - 2363157-2724\\
\small Análisis y Diseño de Algoritmos}
\date{}

\begin{document}
\maketitle

\begin{abstract}
El problema de la \textit{suma de rango} consiste en calcular la suma de los
elementos de un arreglo entre dos índices $[l, r]$.
La solución ingenua recorre todos los elementos del intervalo, logrando
complejidad $O(n)$ por consulta.
Mediante la técnica de \textbf{Prefix Sum}, precomputamos en $O(n)$ un arreglo
auxiliar que permite responder cualquier consulta en tiempo constante $O(1)$.
Las pruebas empíricas muestran que la versión optimizada es hasta
\textbf{30.903 veces más rápida} para $n = 10^6$.
\end{abstract}

\section{Introducción}
Sea $A = [a_0, a_1, \ldots, a_{n-1}]$ un arreglo de $n$ enteros.
La consulta de suma de rango se define como:
\[
\text{query}(l, r) = \sum_{i=l}^{r} a_i, \quad 0 \le l \le r < n
\]
Cuando el número de consultas $q \sim n$, la solución ingenua cuesta
$O(n \cdot q)$, lo que para $n = q = 10^6$ equivale a $10^{12}$ operaciones,
completamente inviable en tiempo real.

\section{Análisis de Complejidad}

\subsection{Implementación Ingenua --- $O(n)$ por consulta}
El enfoque directo itera sobre todos los índices del intervalo $[l, r]$.
Para $q$ consultas la complejidad total es $O(n \cdot q)$.

\subsection{Prefix Sum --- $O(1)$ por consulta}
\textbf{Definición.} Sea $P$ el arreglo de sumas prefijas:
\[
P[0] = 0, \qquad P[i] = \sum_{k=0}^{i-1} a_k, \quad \forall i \in [1, n]
\]

\textbf{Teorema.} Para todo $0 \le l \le r < n$:
\[
\text{query}(l, r) = P[r+1] - P[l]
\]

\textit{Demostración.}
\begin{align*}
P[r+1] - P[l]
&= \sum_{k=0}^{r} a_k - \sum_{k=0}^{l-1} a_k \\
&= \sum_{k=l}^{r} a_k \quad \blacksquare
\end{align*}
Costo total: $O(n)$ de precomputación y $O(1)$ por consulta, resultando en un costo global de $O(n + q)$.

\section{Resultados}
Las pruebas empíricas con 5.000 consultas por prueba muestran:

\begin{table}[h]
\centering
\small
\begin{tabular}{rrrr}
\toprule
$n$ & Ingenua (ms) & Prefix Sum (ms) & Speedup \\
\midrule
1.000      &     19,8  &   0,28 &      70$\times$ \\
10.000     &    196,6  &   0,45 &     440$\times$ \\
100.000    &  2.014,3  &   0,43 &   4.733$\times$ \\
500.000    & 10.265,4  &   0,51 &  19.979$\times$ \\
1.000.000  & 20.463,1  &   0,66 &  30.903$\times$ \\
\bottomrule
\end{tabular}
\caption{Tiempo total de 5.000 consultas por tamaño}
\end{table}

\section{Aplicaciones Reales}
\textbf{Bases de datos.} PostgreSQL usa sumas prefijas en funciones de
ventana (\texttt{SUM OVER}) para calcular rangos en $O(1)$.

\textbf{Visión por Computador.} La Integral Image de Viola-Jones (2001)
aplica prefix sum 2D para detectar rostros en tiempo real.

\textbf{Business Intelligence.} Google Analytics y Apache Druid calculan
métricas acumuladas sobre ventanas de tiempo en $O(1)$.

\begin{thebibliography}{9}
\bibitem{clrs}
T.~Cormen et al., \textit{Introduction to Algorithms}, 4th ed. MIT Press, 2022.
\bibitem{violajones}
P.~Viola, M.~Jones, ``Rapid object detection using a boosted cascade,''
\textit{CVPR}, 2001.
\end{thebibliography}

\vspace{0.5em}
\noindent\textbf{Repositorio:} \\
\small\url{https://github.com/ashlyalex13/Prefix-sum-algoritmos}

\end{document}